\clearpage\fancyhead[L]{\small PALLAS.}\hypertarget{native-361-lower}{}\label{native-361-lower}
\begin{center}\textbf{PALLAS.}\par\vspace{4pt}\hrule\vspace{4pt}\fontsize{10}{12}\selectfont Göttingische gelehrte Anzeigen. No. 127, pp. 1257\;..\;1259, 8 Aug. 1812.\par\vspace{4pt}\hrule\vspace{4pt}\end{center}

\par\indent Issue 67 of this journal [p. 349 of this volume] has already reported the first sighting of Pallas this year and the striking agreement of its position with the ephemeris calculated from the theory of perturbations. Few observations have been made at the local observatory since then. Around the time of opposition in particular, poor weather, moonlight and the planet’s inconveniently high position for this site allowed only one observation. The following are all the observations Professor \textsc{Gauss} made at our observatory this year:

\begin{center}
\fontsize{10}{12}\selectfont\setlength{\tabcolsep}{3pt}\renewcommand{\arraystretch}{1.15}\begin{tabular}{c|c|c|c}
1812 & Mean time & \shortstack{Right\\ascension} & \shortstack{Northern\\declination} \\
\hline
April 9 & $11^{\mathrm u}\;32^{\mathrm m}\;27^{\mathrm s}$ & $269^\circ\;5'\;58''4$ & $16^\circ\;38'\;0''$ \\
May 2 & 10 \quad 36 \quad 6 & 268 \quad 42 \quad 5.3 & 21 \quad 4 \quad 31.7 \\
— 4 & 10 \quad 44 \quad 44 & 268 \quad 31 \quad 38.0 & 21 \quad 25 \quad 4.3 \\
June 7 & 11 \quad 13 \quad 51 & 262 \quad 50 \quad 51.1 & 25 \quad 2 \quad 56.6 \\
\end{tabular}
\end{center}

\par\indent At the last observation, the planet was of the tenth magnitude.

\par\indent Professor \textsc{Gauss} had invited several of his astronomer friends elsewhere, who have better equipment at their disposal, to observe Pallas around the time of opposition. According to the reports received so far, however, they were unable to do so at the Imperial Observatory in Paris or in Marseille or Berlin. This is hardly surprising, since the planet is so faint this year that it is scarcely visible with fixed instruments when even the slightest illumination is used. Only Mr. \textsc{Burckhardt} was fortunate enough to obtain the following meridian observations at the observatory of the military school in Paris:

\begin{center}
\fontsize{10}{12}\selectfont\setlength{\tabcolsep}{3pt}\renewcommand{\arraystretch}{1.15}\begin{tabular}{c|c|c|c}
1812 & Mean time & \shortstack{Right\\ascension} & \shortstack{Northern\\declination} \\
\hline
June 5 & $12^{\mathrm u}\;35^{\mathrm m}\;57^{\mathrm s}4$ & $263^\circ\;15'\;24''0$ & $24^\circ\;58'\;28''5$ \\
6 & 12 \quad 33 \quad 10.5 & 263 \quad 2 \quad 35 & 25 \quad 0 \quad 30 \\
8 & 12 \quad 21 \quad 37.0 & 262 \quad 37 \quad 6.5 & 25 \quad 3 \quad 57.7 \\
11 & 12 \quad 7 \quad 15.7 & 261 \quad 58 \quad 37.8 & 25 \quad 6 \quad 29.3 \\
\end{tabular}
\end{center}

\par\indent Professor \textsc{Gauss} gave these observations to Mr. \textsc{Nicolai} to calculate the opposition from them. His result is as follows:
\SourceLeaf{362}\begin{center}\fontsize{10}{12}\selectfont\setlength{\tabcolsep}{0pt}\renewcommand{\arraystretch}{1.15}\begin{tabular}{l}
Opposition of Pallas 1812 June 10. $3^{\mathrm u}\;2^{\mathrm m}\;48^{\mathrm s}$\\mean time at the Göttingen meridian
\end{tabular}\end{center}
\begin{center}\fontsize{10}{12}\selectfont\setlength{\tabcolsep}{2pt}\renewcommand{\arraystretch}{1.15}\begin{tabular}{lcr}
True longitude & $.\;.\;.\;.\;.\;.\;.\;.\;.\;.\;.\;.\;.\;.\;.\;.\;.\;.$ & $259^\circ\;37'\;51''7$\\
True geocentric latitude & $.\;.\;.\;.\;.\;.\;.\;.\;.\;.\;.\;.\;.$ & $48\;16\;26.1\;\mathrm{N}.$\\
\end{tabular}\end{center}

\par\indent This is now the \emph{eighth} observed opposition of Pallas since its discovery. The result differs from the advance calculation Mr. \textsc{Nicolai} made a year earlier by 29 minutes in time, 1 minute in longitude, and 10 seconds in latitude (see Gött. gel. Anz. 1811, p. 1292 [p. 329 of this volume]). Professor \textsc{Gauss} believes that even in these small differences he can already recognize with certainty the influence of the still unexamined perturbations due to Saturn and Mars. All the preparations for calculating Saturn’s perturbations have already been made, and we will give a fuller account of them in due course. In the meantime Professor \textsc{Gauss} has investigated how closely all eight oppositions observed so far can still be represented if only the perturbations due to Jupiter are taken into account. The result is that the following adjustments are needed to obtain the closest possible agreement.

\begin{center}\fontsize{9}{11}\selectfont\setlength{\tabcolsep}{2pt}\renewcommand{\arraystretch}{1.15}\begin{tabular}{lcr}
\shortstack[l]{the epoch of mean longitude\\on 10 June 1812 by} & $.\;.\;.\;.$ & $-26''95$\\
the daily mean motion by & $.\;.\;.\;.\;.\;.\;.\;.\;.\;.$ & $-0''01225$\\
the logarithm of the semiaxis by & $.\;.\;.\;.\;.\;.\;.\;.\;.\;.\;.$ & $+0.0000046$\\
the longitude of perihelion by & $.\;.\;.\;.\;.\;.\;.\;.\;.\;.\;.\;.\;.\;.$ & $-18''84$\\
the angle of eccentricity by & $.\;.\;.\;.\;.\;.\;.\;.\;.\;.\;.\;.\;.$ & $-3''92$\\
the longitude of the node by & $.\;.\;.\;.\;.\;.\;.\;.\;.\;.\;.\;.\;.\;.$ & $+0''80$\\
the inclination of the orbit by & $.\;.\;.\;.\;.\;.\;.\;.\;.\;.\;.\;.\;.\;.$ & $-1''21$\\
\end{tabular}\end{center}

\Needspace{16\baselineskip}\noindent The observed and calculated opposition then agree as follows:

\begin{center}
\fontsize{10}{12}\selectfont\setlength{\tabcolsep}{7pt}\renewcommand{\arraystretch}{1.15}\begin{tabular}{c|c|c}
 & \multicolumn{2}{c}{Difference in} \\
Opposition & mean longitude & heliocentr. latitude \\
\hline
1803 & $+$\ 6$''$4  & $-$\ 2$''$8 \\
1804 & $-$10.9 & $-$9.7 \\
1805 & $-$13.8 & $-$0.7 \\
1807 & $+$10.2 & $-$5.8 \\
1808 & $+$19.1 & $+$11.7 \\
1809 & $+$16.1 & $-$1.5 \\
1811 & $+$14.2 & $-$9.5 \\
1812 & $-$22.3 & $+$6.1 \\
\end{tabular}
\end{center}

\par\indent According to Mr. \textsc{Nicolai}’s calculation, the \emph{ninth} opposition will occur on 18 August 1813 at $9^{\mathrm u}$, at longitude $325^\circ\;24'$ and geocentric latitude $24^\circ\;37'$~N.
